% Generated draft: requires mathematical and source review.
\documentclass[11pt]{article}
\usepackage[margin=25mm]{geometry}
\usepackage{fontspec}
\setmainfont{Latin Modern Roman}
\usepackage{amsmath,amssymb,mathtools,unicode-math}
\setmathfont{Latin Modern Math}
\usepackage{enumitem,needspace,xurl,hyperref}
\setlength{\parindent}{0pt}
\setlength{\parskip}{5pt}
\setlength{\emergencystretch}{3em}
\newcommand{\N}{\mathbb{N}}
\newcommand{\Z}{\mathbb{Z}}
\newcommand{\Q}{\mathbb{Q}}
\newcommand{\R}{\mathbb{R}}
\newcommand{\C}{\mathbb{C}}
\newcommand{\F}{\mathbb{F}}
\newcommand{\A}{\mathbb{A}}
\newcommand{\PP}{\mathbb{P}}
\newcommand{\E}{\mathbb{E}}
\newcommand{\eps}{\varepsilon}
\newcommand{\dd}{\,\mathrm{d}}
\begin{document}
\section*{UnsolvedMath ID: 800018}
\textbf{Source problem code:} catalogue entry
\subsection{Definitions and mathematical statement}
A commutative semigroup is a nonempty set \(S\) equipped with an associative operation \(\cdot\) satisfying \(x\cdot y=y\cdot x\) for all \(x,y\in S\). No identity element or finite generation hypothesis is imposed. An endomorphism of \(S\) means a map \(f:S\to S\) such that
\[
f(x\cdot y)=f(x)\cdot f(y)\qquad(x,y\in S).
\]
Here endomorphisms are ordinary semigroup homomorphisms, with no additional identity-preservation requirement.

\textbf{Problem 18 (A.A. Borisov).} Describe, up to semigroup isomorphism, all commutative semigroups \(S\) for which
\[
\forall f\in\operatorname{End}(S)\;\forall x,y\in S,
\qquad f(x)=f(y)\Longrightarrow x=y.
\]
Thus the requested description must give necessary and sufficient conditions for every endomorphism to be injective.

\textbf{Supplied finite-generation qualification.} The source record reports that [Borisov4] proves that the finitely generated commutative semigroups with this property are precisely subsemigroups of an infinite monogenic semigroup or finite ordinal sums of such semigroups. An infinite monogenic semigroup is \(\{a,a^2,a^3,\ldots\}\), with distinct positive powers and multiplication \(a^ma^n=a^{m+n}\). The record does not define its ordinal-sum construction or supply the bibliographic identity of [Borisov4]. The treatment of the one-element semigroup also requires clarification: it satisfies the stated endomorphism condition, but its inclusion in the reported classification is not explained. Consequently this classification is retained as a source-qualified report, not asserted here as a verified theorem or used to restrict the general problem. The nonempty-set convention above is explicit; the record does not state whether empty semigroups are admitted.
\subsection{Short English statement}
Classify all commutative semigroups whose every multiplication-preserving self-map is injective. The question is attributed to A.A. Borisov. The supplied record labels it open and reports a finitely generated classification, but that report has unresolved terminology, a missing reference, and a degenerate-case qualification requiring review.
\subsection{Sources}
\begin{itemize}
\item Local UnsolvedMath registry, permanent ID \(800018\), exact source code \texttt{catalogue entry}, title \emph{The Lyapin's Notebook --- Problem 18}. The supplied registry record labels the problem \texttt{open}; this is supplied curation, not an independently verified status assessment.
\item Bershadsky, Kublanovsky, Mashevitzky, \emph{The Lyapin's Notebook: Unsolved Problems in Semigroup Theory} (2026), Problem 18, \url{https://arxiv.org/abs/2604.04763}. The source record attributes the question itself to A.A. Borisov and lists Bershadsky, Kublanovsky, Mashevitzky in its proposer metadata. It records access on 2026-10-16 and extraction by \texttt{arxiv\_latex}; these are supplied provenance annotations, not a claim of independent source access.
\item [Borisov4], cited in the supplied statement for the finitely generated classification. No author expansion, title, publication details, or separate URL for this reference is supplied.
\end{itemize}
\end{document}
